\section{Grobe Flow-Analyse}
\label{sec:cashflow}

Die Aufgabenstellung gibt f\"ur jede Option die Investitionssumme, die
Kapitalbindung und den erwarteten EBITDA-Zuwachs vor. Nicht vorgegeben sind
deren zeitliche Verteilung, die Abgrenzung des ausgewiesenen Cash-Flows, die
steuerliche Betrachtung und der Horizont der Renditerechnung. Diese Punkte
sind Annahmen dieser Arbeit und werden vorab offengelegt; alle folgenden
Tabellen sind aus ihnen berechnet.

\subsection*{Annahmen der Flow-Analyse}
\addcontentsline{toc}{subsection}{Annahmen der Flow-Analyse}

\annahme{\textbf{Anfangsbestand.} Die Rechnung setzt auf der
Nettofinanz\-liquidit\"at von \MioEUR{\NettoliquiditaetZFIV} zum
31.~Dezember 2024 auf, wie in der Aufgabenstellung vorgegeben. Zum
31.~Dezember 2025 wies die Gesellschaft tats\"achlich
\MioEUR{\NettoliquiditaetZFV} aus\quelleGB{2025}{4}; der Bestand ist trotz
Dividende und Aktienr\"uckkauf gestiegen. Die Vorgabe der Aufgabenstellung
wird gleichwohl beibehalten, damit die Ergebnisse vergleichbar bleiben.}

\annahme{\textbf{Investitionsprofil.} Die Investitionen der Optionen~A und~B
verteilen sich gleichm\"a{\ss}ig auf die Jahre 2026 bis 2028. Bei Option~C
wird der Kaufpreis im Vollzugsjahr 2026 f\"allig; nur die
Integrationsaufwendungen verteilen sich auf 2027 und 2028. Ein Kaufpreis wird
nicht in Dritteln gezahlt, weshalb hier bewusst von der gleichm\"a{\ss}igen
Verteilung abgewichen wird.}

\annahme{\textbf{Kapitalbindung.} Der Aufbau der Kapitalbindung erfolgt in dem
Jahr beziehungsweise den Jahren unmittelbar vor dem jeweiligen
Ergebnisbeitrag, da das Umlaufverm\"ogen den Umsatz vorfinanziert. Bei
Option~C f\"allt das Umlaufverm\"ogen des erworbenen Gesch\"afts mit dem
Vollzug an.}

\annahme{\textbf{Abgrenzung des operativen Cash-Flows.} Ausgewiesen wird der
\emph{zus\"atzliche} operative Cash-Flow der jeweiligen Option, das hei{\ss}t
der EBITDA-Zuwachs abz\"uglich des Aufbaus der Kapitalbindung. Der operative
Cash-Flow des bestehenden Kerngesch\"afts bleibt unber\"ucksichtigt, damit die
Wirkung der Entscheidung sichtbar bleibt und nicht im Konzern-Cash-Flow
untergeht. Die ausgewiesene Nettoliquidit\"at ist entsprechend keine Prognose
des Konzernbestands, sondern der Bestand unter alleiniger Wirkung der Option.}

\annahme{\textbf{Steuern und Abschreibungen.} Die Rechnung ist eine
Vorsteuerbetrachtung. Der vorgegebene EBITDA-Zuwachs geht ungek\"urzt in den
operativen Cash-Flow ein; Ertragsteuern und die Abschreibungen auf die
Investitionssummen bleiben unber\"ucksichtigt, weil die Aufgabenstellung die
Ergebnissteigerung ausdr\"ucklich als EBITDA-Gr\"o{\ss}e vorgibt und weder
einen Steuersatz noch eine Nutzungsdauer nennt. Die ausgewiesenen internen
Zinsf\"u{\ss}e sind deshalb Vorsteuerrenditen und liegen \"uber den
entsprechenden Nachsteuerwerten. Da alle drei Optionen gleich behandelt
werden, bleibt die Rangfolge des Vergleichs davon unber\"uhrt; die absoluten
Renditeangaben sind entsprechend zu lesen.}

\annahme{\textbf{Renditehorizont.} Der interne Zinsfu{\ss} wird \"uber den
Horizont 2026 bis 2035 ermittelt, also zehn Jahre ab dem ersten
Investitionsjahr. Zehn Jahre sind ein \"ublicher Konventionswert f\"ur die
mittelfristige Investitionsrechnung und liegen deutlich \"uber der
dreij\"ahrigen Investitionsphase aller drei Optionen, sodass sich auch der
sp\"ate Ergebnisbeitrag von Option~A noch \"uber mehrere Jahre auswirken
kann. Weil diese Wahl eine Konvention und keine aus den Unternehmensdaten
abgeleitete Gr\"o{\ss}e ist, wird sie nicht isoliert verwendet: Der
EBITDA-Zuwachs l\"auft ab seinem Startjahr konstant weiter, die
Kapitalbindung wird im letzten Jahr des Horizonts freigesetzt, und
Tabelle~\ref{tab:irr-sens} zeigt denselben Vergleich zus\"atzlich f\"ur
f\"unf und f\"unfzehn Jahre. Ein interner Zinsfu{\ss} ohne erkl\"arten
Horizont ist nicht aussagef\"ahig.}

\annahme{\textbf{Verh\"altnis zur eigenen Sch\"atzung aus
Abschnitt~\ref{sec:optionen}.} Investitionssumme und Kapitalbindung sind
f\"ur diesen Abschnitt durch die Aufgabenstellung vorgegeben und weichen
von der eigenen Sch\"atzung aus Abschnitt~\ref{sec:optionen} ab, die auf
der aktivierten Entwicklungsquote und der Nettofinanzliquidit\"at
aufbaut: \MioEUR{\InvestitionssummeA} statt \MioEUR{\InvestitionsbedarfEigenA}
bei Option~A, \MioEUR{\InvestitionssummeB} statt
\MioEUR{\InvestitionsbedarfEigenB} bei Option~B und
\MioEUR{\InvestitionssummeC} statt \MioEUR{\InvestitionsbedarfEigenC} bei
Option~C. Beide Rechnungen bleiben nebeneinander bestehen; die folgenden
Tabellen folgen durchg\"angig der Vorgabe dieses Abschnitts.}

\paragraph{Vorbehalt zur EBITDA-Marge Ende 2028.}
Die in Tabelle~\ref{tab:vergleich2028} ausgewiesene EBITDA-Marge bezieht den
EBITDA-Zuwachs der Option auf das operative EBITDA und den Umsatz des
Gesch\"aftsjahres 2025, die konstant fortgeschrieben werden. Diese Basis ist
nicht neutral: Der Umsatz 2025 lag mit \MioEUR{\UmsatzZFV} um
\PctBetrag{\UmsatzDeltaZFV} unter dem Umsatz 2024 von
\MioEUR{\UmsatzZFIV}.\quelleGB{2025}{2} Der Nenner ist damit klein, und die
ausgewiesene Marge f\"allt entsprechend hoch aus. Sie zeigt damit nicht,
wie profitabel die Option wirklich ist; sie spiegelt zu einem erheblichen
Teil den Umsatzeinbruch des Basisjahres. Tabelle~\ref{tab:marge-sens} stellt dem
Basisfall daher zwei gr\"o{\ss}ere Nenner gegen\"uber: ein Umsatzwachstum von
\Pct{\MargenWachstumSzenario} pro Jahr \"uber die drei Jahre bis 2028 sowie
eine R\"uckkehr auf das Umsatzniveau des Jahres 2024. Die Rechnung
unterstellt nicht, dass eines dieser Niveaus erreicht wird; sie zeigt allein,
wie stark die Kennzahl von ihrem Nenner abh\"angt.

\begin{table}[!htbp]
\centering
\caption{Sensitivit\"at der EBITDA-Marge Ende 2028 gegen\"uber der
angesetzten Umsatzbasis.}
\label{tab:marge-sens}
\begin{tabularx}{\linewidth}{@{}Xrrr@{}}
\toprule
\textbf{Umsatzbasis 2028} & \textbf{Option A} & \textbf{Option B} & \textbf{Option C}\\
\midrule
\tabellenkoerper{data/marge_sensitivitaet}
\bottomrule
\end{tabularx}
\end{table}
\FloatBarrier

\subsection{Cash-Flow-Tabellen der Optionen}
\label{subsec:cf-tabellen}

Alle Betr\"age in Mio.\,\euro. Die Spalte \enquote{Operativer CF} weist den
zus\"atzlichen operativen Cash-Flow der Option aus, die Spalte
\enquote{Nettoliquidit\"at} den Bestand nach Investition und Kapitalbindung.

\begin{table}[!htbp]
\centering
\caption{Option~A: Diversifikation \enquote{Next Automation}.}
\label{tab:cf-a}
\small\begin{tabularx}{\linewidth}{@{}lrrrX@{}}
\toprule
\textbf{Jahr} & \textbf{Investition} & \textbf{Operativer CF}
& \textbf{Netto\-liquidit\"at nach Investition} & \textbf{Bemerkung}\\
\midrule
\tabellenkoerper{data/cf_a}
\bottomrule
\end{tabularx}
\end{table}

\begin{table}[!htbp]
\centering
\caption{Option~B: Batterie- und Brennstoffzellen-Produktionstechnik.}
\label{tab:cf-b}
\small\begin{tabularx}{\linewidth}{@{}lrrrX@{}}
\toprule
\textbf{Jahr} & \textbf{Investition} & \textbf{Operativer CF}
& \textbf{Netto\-liquidit\"at nach Investition} & \textbf{Bemerkung}\\
\midrule
\tabellenkoerper{data/cf_b}
\bottomrule
\end{tabularx}
\end{table}

\begin{table}[!htbp]
\centering
\caption{Option~C: Anorganisches Wachstum durch gezielte Akquisition.}
\label{tab:cf-c}
\small\begin{tabularx}{\linewidth}{@{}lrrrX@{}}
\toprule
\textbf{Jahr} & \textbf{Investition} & \textbf{Operativer CF}
& \textbf{Netto\-liquidit\"at nach Investition} & \textbf{Bemerkung}\\
\midrule
\tabellenkoerper{data/cf_c}
\bottomrule
\end{tabularx}
\end{table}
\FloatBarrier

% Auf 6.2 folgt unmittelbar die Unterunterueberschrift 6.2.1; ohne
% \needspace standen beide Ueberschriften ohne jeden Rumpf am Fuss von
% Seite 22. Bemessen fuer beide Ueberschriften und die ersten Textzeilen.
\needspace{8\baselineskip}
\subsection{Ergebnis der einzelnen Flow-Analysen}
\label{subsec:cf-ergebnis}

\subsubsection{Option A}
% FAKTENBASIS 6.2 A
%  - Kapitalbedarf: 35 Mio. EUR Investition + 15 Mio. EUR Kapitalbindung
%  - EBITDA-Zuwachs +16 Mio. EUR erst ab 2028 -> spaeteste Ergebniswirkung,
%    aber groesster absoluter Zuwachs der drei Optionen
%  - Liquiditaetsverlauf und Endbestand: Tabelle \ref{tab:cf-a}
%    (126,5 -> 107,3 -> 104,2; kein Jahr mit negativem Bestand)
%  - Erwartete Rendite: Tabelle \ref{tab:vergleich2028}; auf fuenf Jahre
%    deutlich niedriger, siehe Tabelle \ref{tab:irr-sens}
%  - Relation Kapitalbedarf / Liquiditaetswirkung: rund ein Viertel der
%    Ausgangsliquiditaet \NettoliquiditaetZFIV Mio. EUR wird aufgezehrt
\bk{Option~A verlangt den zweith\"ochsten Kapitaleinsatz des Vergleichs:
\MioEUR{\InvestitionssummeA} Investition und \MioEUR{\KapitalbindungssummeA}
Kapitalbindung, zusammen \MioEUR{\KapitalbedarfA}. Die erwartete Rendite
liegt bei \Pct{\IrrA} \"uber zehn Jahre, \"uber f\"unf Jahre jedoch nur bei
\Pct{\IrrFuenfA}; das ist eine Folge davon, dass der Ergebnisbeitrag von
\MioEUR{\EbitdaZuwachsA} erst \EbitdaStartA\ einsetzt. In Relation dazu ist
die Liquidit\"atswirkung unauff\"allig: Der Bestand sinkt gleichm\"a{\ss}ig
auf \MioEUR{\LiquiditaetEndeA} und unterschreitet in keinem Jahr eine
kritische Schwelle. Der gesamte Mittelabfluss liegt jedoch vor dem ersten
R\"uckfluss; die Option ist finanzierbar, aber nicht selbsttragend.}

\subsubsection{Option B}
% FAKTENBASIS 6.2 B
%  - Kapitalbedarf: 30 Mio. EUR Investition + 15 Mio. EUR Kapitalbindung
%    -> geringster Kapitalbedarf der drei Optionen
%  - EBITDA-Zuwachs +14 Mio. EUR bereits ab 2027
%  - Liquiditaetsverlauf und Endbestand: Tabelle \ref{tab:cf-b}
%    (120,7 -> 117,2 -> 121,2; Bestand 2028 wieder ueber dem Vorjahr)
%  - Erwartete Rendite: hoechster IRR aller Horizonte,
%    Tabellen \ref{tab:vergleich2028} und \ref{tab:irr-sens}
%  - Relation Kapitalbedarf / Liquiditaetswirkung: der Ergebnisbeitrag
%    setzt ein, bevor das Investitionsprogramm auslaeuft
\bk{Option~B hat mit \MioEUR{\KapitalbedarfB} den geringsten Kapitalbedarf
(\MioEUR{\InvestitionssummeB} Investition und
\MioEUR{\KapitalbindungssummeB} Kapitalbindung) und zugleich die h\"ochste
erwartete Rendite: \Pct{\IrrB} \"uber zehn und bereits \Pct{\IrrFuenfB}
\"uber f\"unf Jahre. Der Grund liegt im fr\"uhen Ergebnisbeitrag ab
\EbitdaStartB. Entsprechend g\"unstig f\"allt die Liquidit\"atswirkung aus:
Der Bestand erreicht seinen Tiefpunkt mit \MioEUR{\LiquiditaetTiefB} bereits
2027 und liegt Ende 2028 mit \MioEUR{\LiquiditaetEndeB} wieder h\"oher. Als
einzige der drei Optionen tr\"agt sie ihre letzte Investitionsrate aus dem
eigenen Ergebnisbeitrag. Kapitalbedarf, Rendite und Liquidit\"atswirkung
weisen damit in dieselbe Richtung.}

\subsubsection{Option C}
% FAKTENBASIS 6.2 C
%  - Kapitalbedarf: 60 Mio. EUR Investition + 20 Mio. EUR Kapitalbindung,
%    davon 65 Mio. EUR bereits 2026 zahlungswirksam
%  - EBITDA-Zuwachs +10 Mio. EUR ab 2027, geringster absoluter Zuwachs
%  - Liquiditaetsverlauf und Endbestand: Tabelle \ref{tab:cf-c}
%    (Tiefpunkt 73,2 in 2026, danach 75,7 und 78,2)
%  - Erwartete Rendite: niedrigster IRR; ueber fuenf Jahre negativ,
%    Tabelle \ref{tab:irr-sens}
%  - Relation Kapitalbedarf / Liquiditaetswirkung: knapp die Haelfte der
%    Ausgangsliquiditaet \NettoliquiditaetZFIV Mio. EUR ist in einem Jahr
%    gebunden; Net Debt/EBITDA bleibt negativ, also Nettoguthaben
%  - Der Ergebnisbeitrag wird gekauft, nicht aufgebaut -> Risiko liegt in
%    Kaufpreis und Integration, nicht in der Markterschliessung
\bk{Option~C bindet mit \MioEUR{\KapitalbedarfC} den h\"ochsten Betrag und
liefert mit \MioEUR{\EbitdaZuwachsC} den geringsten Ergebnisbeitrag. Die
erwartete Rendite bleibt entsprechend niedrig: \Pct{\IrrC} \"uber zehn Jahre
und \Pct{\IrrFuenfC} \"uber f\"unf, also negativ. Die Liquidit\"atswirkung
ist zudem die ung\"unstigste, weil sie sich nicht verteilt:
\MioEUR{\MittelabflussErstesJahrC} flie{\ss}en allein 2026 ab, der Bestand
f\"allt auf \MioEUR{\LiquiditaetTiefC} und erreicht Ende 2028 lediglich
\MioEUR{\LiquiditaetEndeC}. Der Kapitalbedarf steht damit in keinem
g\"unstigen Verh\"altnis zu Rendite und Liquidit\"atswirkung; die Option
bleibt tragbar, verbraucht aber den finanziellen Spielraum f\"ur mehrere
Jahre.}

\subsection{Vergleich der Ergebniswirkung (Ende 2028)}
\label{subsec:vergleich}

Tabelle~\ref{tab:vergleich2028} stellt die Ergebniswirkung der drei
Optionen zum Jahresende 2028 gegen\"uber. Da der interne Zinsfu{\ss} ohne
Angabe des unterstellten Horizonts nicht aussagef\"ahig ist, zeigt
Tabelle~\ref{tab:irr-sens} denselben Vergleich zus\"atzlich f\"ur drei
Horizontl\"angen.

\begin{table}[!htbp]
\centering
\caption{Vergleich der drei Optionen zum Jahresende 2028. Die EBITDA-Marge
ist auf die konstant fortgeschriebene Umsatzbasis 2025 bezogen; zum
Vorbehalt siehe Tabelle~\ref{tab:marge-sens}.}
\label{tab:vergleich2028}
\begin{tabularx}{\linewidth}{@{}Xrrr@{}}
\toprule
\textbf{Kennzahl} & \textbf{Option A} & \textbf{Option B} & \textbf{Option C}\\
\midrule
\tabellenkoerper{data/vergleich2028}
\bottomrule
\end{tabularx}
\end{table}

\begin{table}[!htbp]
\centering
\caption{Sensitivit\"at des internen Zinsfu{\ss}es gegen\"uber der
Horizontl\"ange.}
\label{tab:irr-sens}
\begin{tabularx}{\linewidth}{@{}Xrrr@{}}
\toprule
\textbf{Horizont} & \textbf{Option A} & \textbf{Option B} & \textbf{Option C}\\
\midrule
\tabellenkoerper{data/irr_sensitivitaet}
\bottomrule
\end{tabularx}
\end{table}
\FloatBarrier

\subsubsection{Fazit des Vergleichs}
% FAKTENBASIS 6.3
%  - Kumulierte Investition: A 35, B 30, C 60 Mio. EUR;
%    Kapitalbindung A 15, B 15, C 20 Mio. EUR
%  - EBITDA-Zuwachs: A +16 ab 2028, B +14 ab 2027, C +10 ab 2027
%  - EBITDA-Marge Ende 2028, IRR, Liquiditaetsreserve:
%    Tabelle \ref{tab:vergleich2028}
%  - IRR-Sensitivitaet nach Horizontlaenge: Tabelle \ref{tab:irr-sens};
%    Rangfolge B > A > C in allen drei Horizonten, C auf fuenf Jahre negativ
%  - Margen-Sensitivitaet: Tabelle \ref{tab:marge-sens}; die Rangfolge
%    A > B > C bleibt in allen drei Szenarien erhalten, das Niveau nicht
%  - Net Debt/EBITDA nur fuer Option C relevant, da A und B die
%    Nettoliquiditaet nicht nennenswert aufzehren; der Wert bleibt negativ
%  - Umsatzbasis der Margenrechnung: Geschaeftsjahr 2025, konstant
%    fortgeschrieben; Umsatzrueckgang 2025 \UmsatzDeltaZFV %
%  - Die drei Optionen sind in der Rechnung unabhaengig; die
%    Ausgangsliquiditaet \NettoliquiditaetZFIV Mio. EUR traegt A und B
%    zusammen, nicht alle drei
\bk{Die Rangfolge der drei Optionen ist \"uber alle betrachteten Kennzahlen
hinweg stabil, und sie f\"allt nicht zugunsten des gr\"o{\ss}ten
Kapitaleinsatzes aus. Beim internen Zinsfu{\ss} f\"uhrt Option~B mit
\Pct{\IrrB} vor Option~A mit \Pct{\IrrA} und Option~C mit \Pct{\IrrC}; diese
Reihenfolge gilt in allen drei Horizonten der Tabelle~\ref{tab:irr-sens}.
Bemerkenswert ist der Abstand zwischen den Horizonten: \"Uber f\"unf Jahre
erreicht Option~A nur \Pct{\IrrFuenfA} und Option~C mit \Pct{\IrrFuenfC}
einen negativen Wert. Eine Renditeaussage ohne Angabe des Horizonts w\"are
hier also nicht nur ungenau, sondern irref\"uhrend.

Beim Kapitalbedarf und bei der Liquidit\"at zeigt sich derselbe Befund. Die
Optionen~A und~B lassen zum Jahresende 2028 Reserven von
\MioEUR{\LiquiditaetEndeA} beziehungsweise \MioEUR{\LiquiditaetEndeB} und
zehren die Nettoliquidit\"at damit nicht auf; f\"ur sie ist die Kennzahl Net
Debt/EBITDA ohne Aussagekraft. Nur bei Option~C ist sie \"uberhaupt
sinnvoll ausweisbar, und auch dort bleibt sie mit \num{\NetDebtEbitdaC}
negativ. Es besteht also weiterhin ein Nettoguthaben. Die Belastung der
Option~C liegt nicht in einer Verschuldung, sondern in der Konzentration des
Mittelabflusses auf ein einziges Jahr und in dem daraus folgenden Verlust an
Handlungsf\"ahigkeit.

Bei der EBITDA-Marge Ende 2028 ist Zur\"uckhaltung geboten. Der Basisfall
weist f\"ur Option~A \Pct{\MargeZweitausendachtundzwanzigA}, f\"ur Option~B
\Pct{\MargeZweitausendachtundzwanzigB} und f\"ur Option~C
\Pct{\MargeZweitausendachtundzwanzigC} aus. Diese Werte beruhen auf dem
konstant fortgeschriebenen Umsatz des Gesch\"aftsjahres 2025, der gegen\"uber
2024 um \PctBetrag{\UmsatzDeltaZFV} niedriger lag. Sie messen damit zu einem
erheblichen Teil den Umsatzeinbruch des Basisjahres und nicht, wie
profitabel die Option wirklich ist. Tabelle~\ref{tab:marge-sens} zeigt, wie stark der
Effekt ist: Bei einer R\"uckkehr auf das Umsatzniveau 2024 f\"allt die Marge
jeder Option um mehrere Prozentpunkte. Die Rangfolge zwischen den Optionen
bleibt in allen Szenarien erhalten, das Niveau nicht: Die Marge taugt hier
zum Vergleich der Optionen untereinander, nicht als absolute Prognose.

Als Fazit ist Option~B die finanzwirtschaftlich \"uberlegene
Einzelentscheidung: geringster Einsatz, h\"ochste Rendite, fr\"uhester
R\"uckfluss. Option~A ist teurer und langsamer, liefert aber den
gr\"o{\ss}ten absoluten Ergebnisbeitrag und ist die einzige Option, die die
Abh\"angigkeit vom Investitionszyklus der Automobilindustrie selbst
verringert; dieser Nutzen erscheint in keiner der Kennzahlen dieser Tabellen.
Option~C ist unter den gegebenen Annahmen die schw\"achste Verwendung der
Mittel. Die Ausgangsliquidit\"at von \MioEUR{\NettoliquiditaetZFIV} tr\"agt
die Optionen~A und~B gemeinsam, nicht aber alle drei; die Entscheidung ist
daher eine Auswahl und keine Reihung.}